%  !TeX  root  =  user_guide.tex

\chapter{Использование внешних Python расширений}\label{sec:external_plugins}\index{plugins}

% when the revision of a section has been finalized,
% comment out the following line:
% \updatedisclaimer

Внешние расширения QGIS написаны на Python. Они находятся либо в <<Официальном>>
либо в <<Пользовательском>> репозитории QGIS, или в репозиториях, поддерживаемых
отдельными авторами.
В таблице \ref{tab:external_plugins} перечислены <<официальные>> расширения с
кратким описанием.
Подробная документация, минимальная версия QGIS, домашняя страница, авторы
ии другая важная информация распространяется вместе с внешними расширениями
и не включена в это руководство.
\footnote{Обновления для расширений ядра могут быть поставляться через этот
репозиторий как оверлеи.}
\footnote{fTools, Mapserver Export и Plugin Installer являются Python
расширениями, но они "--- часть исходного кода QGIS, и автоматически
загружаются и активируются в Менеджере расширений QGIS
(смотри Раздел~\ref{sec:load_external_plugin}).}

Актуальный список <<официальных>> расширений доступен на странице Официального
репозитория QGIS \url{http://qgis.osgeo.org/download/plugins.html}. Так же
этот список доступен автоматически в \filename{Установщике расширений}
из меню \dropmenuopttwo{plugin_installer}{Загрузить модули...}.

\begin{table}[H]
\centering
 \begin{tabular}{|l|l|p{8cm}|}
\hline \textbf{Иконка} & \textbf{Название расширения} & \textbf{Описание}\\
\hline
\includegraphics[width=0.7cm]{zoom2point_icon}
 & Zoom To Point \index{plugins!Zoom To Point} & Переход к координатам
  заданным в окне ввода. Можно задавать как масштаб, так и
  контролировать охват.\\
\hline
\includegraphics[width=0.7cm]{plugin_installer}
 & Plugin Installer \index{plugins!Plugin Installer} & Последняя версия Установщика Python расширений.\\
\hline
\end{tabular}
\caption{Официальные внешние расширения QGIS}\label{tab:external_plugins}
\end{table}

Подробное описание процесса установки внешних python расширений доступно в
Разделе \ref{sec:load_external_plugin}.

\begin{Tip} \caption{\textsc{Добавление репозиториев}}
Чтобы добавить <<Пользовательский>> репозиторий и/или несколько авторских
репозиториев, откройте Установщик модулей (\mainmenuopt{Модули} \arrow
\dropmenuopttwo{plugin_installer}{Загрузить модули...}), перейдите на
вкладку \tab{Репозитории} и нажмите \button{Добавить сторонние репозитории}.
Если некоторые из добавленных репозиториев вам не нужны, они могут быть
отключены кнопкой \button{Изменить...}, или полностью удалены нажатием
на кнопку \button{Удалить}.
\end{Tip}
\FloatBarrier
